\chapter{Kurven, die man durchlaufen kann}

\section{Eine Gerade durch den Knoten}

\begin{merksatz}
\boxtitel{Die Kubik mit Schleife, parametrisiert}
$C\colon y^{2} = x^{2}(x + 1)$ hat im Ursprung einen Knoten. Eine Gerade $y = tx$ durch den Knoten schneidet $C$ nach Bézout in drei Punkten; zwei davon sind der Knoten selbst (Vielfachheit $2$). Der dritte:
\[ t^{2}x^{2} = x^{2}(x + 1) \iff x = t^{2} - 1, \qquad y = t(t^{2} - 1). \]

\begin{center}
{\small
\renewcommand{\arraystretch}{1.46}
\begin{tabular}{@{}P{1.0cm}P{1.4cm}P{1.4cm}P{4.6cm}@{}}
\toprule
$t$ & $x$ & $y$ & \textbf{Bemerkung}\\
\midrule
$0$ & $-1$ & $0$ & Scheitel der Schleife\\
$\pm1$ & $0$ & $0$ & der Knoten --- zweimal, aus zwei Richtungen\\
$2$ & $3$ & $6$ & Probe: $36 = 9\cdot4$ \checkmark\\
$\tfrac12$ & $-\tfrac34$ & $-\tfrac38$ & Probe: $\tfrac{9}{64} = \tfrac{9}{16}\cdot\tfrac14$ \checkmark\\
$\infty$ & $\infty$ & $\infty$ & der Punkt im Unendlichen $(0 : 1 : 0)$\\
\bottomrule
\end{tabular}}
\end{center}

\begin{center}
\begin{minipage}{0.9\textwidth}
\textbf{Jeder rationale Wert $t$ gibt einen rationalen Punkt, jeder Punkt außer dem Knoten genau ein $t$: Die Kurve ist ein Bild der Geraden, und ihre rationalen Punkte sind so zahlreich wie die rationalen Zahlen. Die glatte Kubik des Titelbilds erlaubt das nicht --- ihre rationalen Punkte sind eine Gruppe vom Rang $1$, eine Kopie von $\mathbb Z$, nicht von $\mathbb Q$.}
\end{minipage}
\end{center}

\begin{center}
\begin{minipage}{0.88\textwidth}
\kursiv{\textbf{Eine Kurve heißt rational, wenn sie eine Parametrisierung durch rationale Funktionen einer Veränderlichen hat.} Der Kreis ist rational (Kapitel 1), die Neilsche Parabel, die Kubik mit Schleife, jede Gerade und jeder Kegelschnitt mit einem Punkt. Die glatte Kubik ist es nicht, und das ist der Unterschied zwischen Geschlecht null und eins, Kapitel 12. Dieses Kapitel sagt, wie man parametrisiert, wenn es geht --- der Trick ist immer eine Gerade durch einen Punkt, der genügend Schnittpunkte aufbraucht ---, warum es auf dem Grundkörper ankommt, und was Lüroth über alle Parametrisierungen sagt. Für den Leitfall ist die Hyperbel rational, parametrisiert durch $k_{12}$: Darum konnte Kapitel 1 die Lösungsmenge als Tabelle aufschreiben.}
\end{minipage}
\end{center}
\end{merksatz}

\section{Kegelschnitte}

\begin{merksatz}
\boxtitel{Satz 11.1 --- Kegelschnitte mit einem Punkt sind rational}
Sei $C\colon f(x, y) = 0$ ein irreduzibler Kegelschnitt (Grad $2$) über $k$ mit einem Punkt $P_0\in C(k)$. Dann gibt es eine rationale Parametrisierung $t\mapsto(x(t), y(t))$ mit Koeffizienten in $k$, die eine Bijektion zwischen $k\cup\{\infty\}$ und $C(k)$ ist (bis auf höchstens einen Punkt).

\schritt{Beweis}
Die Gerade durch $P_0 = (x_0, y_0)$ mit Steigung $t$: $y = y_0 + t(x - x_0)$. Einsetzen in $f$ gibt ein Polynom zweiten Grades in $x$ mit Koeffizienten in $k[t]$, das die Nullstelle $x_0$ hat (denn $P_0\in C$). Die zweite Nullstelle ist nach Vieta $x(t) = -x_0 - \frac{\text{Koeffizient von } x}{\text{Koeffizient von } x^{2}}$, rational in $t$ mit Koeffizienten in $k$; $y(t)$ aus der Geradengleichung. Umgekehrt bestimmt jeder Punkt $P\ne P_0$ von $C(k)$ die Steigung $t = (y - y_0)/(x - x_0)\in k$ der Geraden $P_0P$ (bzw.\ $t = \infty$ für die Senkrechte). Höchstens die Tangente in $P_0$ gibt keinen zweiten Punkt. $\blacksquare$

\schritt{Beispiel: $x^{2} + y^{2} = 5$ mit $P_0 = (1, 2)$}
$y = 2 + t(x - 1)$: $x^{2} + 4 + 4t(x - 1) + t^{2}(x - 1)^{2} = 5$. Koeffizienten: $x^{2}(1 + t^{2}) + x(4t - 2t^{2}) + (4 - 4t + t^{2} - 5) = 0$. Zweite Nullstelle (Vieta, erste ist $x_0 = 1$): $x = -1 - \frac{4t - 2t^{2}}{1 + t^{2}} = \frac{t^{2} - 4t - 1}{1 + t^{2}}$; damit $x - 1 = \frac{-4t - 2}{1 + t^{2}}$ und $y = 2 + t(x - 1) = \frac{2 - 2t - 2t^{2}}{1 + t^{2}}$. Probe $t = 1$: $x = -2$, $y = -1$: $4 + 1 = 5$. \checkmark\ $t = 0$: $(-1, 2)$. \checkmark

\begin{center}
\begin{minipage}{0.9\textwidth}
\textbf{Der ganze Trick ist Vieta: Eine Gerade trifft einen Kegelschnitt zweimal, ein Schnittpunkt ist bekannt, der andere ist rational in der Steigung. Für eine Kubik durch einen glatten Punkt bleiben zwei unbekannte Schnittpunkte --- deshalb geht es dort nicht, es sei denn, der Punkt ist ein Knoten und verbraucht zwei.}
\end{minipage}
\end{center}

\begin{center}
\begin{minipage}{0.88\textwidth}
\kursiv{\textbf{Ob der Kegelschnitt einen Punkt hat, ist die eigentliche Frage}, und sie hängt vom Körper ab: $x^{2} + y^{2} = 3$ hat keinen rationalen Punkt. Beweis: Gäbe es $x = a/c$, $y = b/c$ in gekürzter Form mit $a^{2} + b^{2} = 3c^{2}$, so wäre $a^{2} + b^{2}\equiv0\pmod3$, und da Quadrate modulo $3$ nur $0$ oder $1$ sind, $a\equiv b\equiv0$; dann $9\mid3c^{2}$, $3\mid c$ --- nicht gekürzt. Also hat der Kreis vom Radius $\sqrt3$ keine rationale Parametrisierung, obwohl er über $\mathbb R$ eine hat und \enquote{dieselbe Kurve} ist wie $x^{2} + y^{2} = 1$. \Person{Legendre} gab 1785 das vollständige Kriterium: Ein Kegelschnitt über $\mathbb Q$ hat genau dann einen rationalen Punkt, wenn er einen reellen und für jede Primzahl $p$ einen $p$-adischen hat (Hasse--Minkowski, Band 1). Für Kubiken ist das falsch: Selmers Kurve $3x^{3} + 4y^{3} + 5z^{3} = 0$ hat Punkte überall lokal und keinen rationalen.}
\end{minipage}
\end{center}
\end{merksatz}

\section{Singuläre Kubiken und höhere Kurven}

\begin{merksatz}
\boxtitel{Satz 11.2 --- Parametrisieren durch den singulären Punkt}
Eine irreduzible ebene Kurve vom Grad $d$ mit einem Punkt der Vielfachheit $d - 1$ ist rational: Die Geraden durch diesen Punkt treffen die Kurve in genau einem weiteren Punkt, und dieser hängt rational von der Steigung ab.

\schritt{Beweis}
Nach Verschieben liegt der Punkt im Ursprung, und $f = f_{d-1} + f_d$ mit homogenen Anteilen der Grade $d - 1$ und $d$ (Terme kleineren Grades fehlen, weil der Ursprung Vielfachheit $d - 1$ hat; höhere gibt es nicht). Auf $y = tx$: $f(x, tx) = x^{d-1}f_{d-1}(1, t) + x^{d}f_d(1, t) = 0$, also $x = 0$ ($(d-1)$-fach) oder $x(t) = -f_{d-1}(1, t)/f_d(1, t)$, $y(t) = t\,x(t)$. $\blacksquare$

\schritt{Beispiele}
\begin{center}
{\small
\renewcommand{\arraystretch}{1.5}
\begin{tabular}{@{}P{3.0cm}P{2.0cm}P{3.6cm}@{}}
\toprule
\textbf{Kurve} & \textbf{Punkt} & \textbf{Parametrisierung}\\
\midrule
$y^{2} = x^{3}$ (Neil) & Spitze, Vielf.\ $2$ & $(t^{2}, t^{3})$\\
$y^{2} = x^{2}(x + 1)$ & Knoten & $(t^{2} - 1,\ t^{3} - t)$\\
$x^{3} + y^{3} = 3xy$ (Folium) & Knoten & $\bigl(\frac{3t}{1 + t^{3}}, \frac{3t^{2}}{1 + t^{3}}\bigr)$\\
$(x^{2} + y^{2})^{2} = x^{2} - y^{2}$ (Lemniskate) & Vielf.\ $2$ --- nicht $3$ & \textbf{trotzdem} rational: $\bigl(\frac{t + t^{3}}{1 + t^{4}}, \frac{t - t^{3}}{1 + t^{4}}\bigr)$\\
Leitfall-Hyperbel $g = 0$ & Kegel\-schnitt & $k_{20} = \tfrac25 - \frac{3}{50k_{12}}$\\
\bottomrule
\end{tabular}}
\end{center}

\begin{center}
\begin{minipage}{0.9\textwidth}
\textbf{Die Lemniskate hat Grad $4$ und im Ursprung nur Vielfachheit $2$; der Satz greift nicht --- und sie ist trotzdem rational, weil sie zwei weitere Doppelpunkte im Unendlichen hat (die Kreispunkte). Das Kriterium ist nicht ein Punkt hoher Vielfachheit, sondern die Gesamtzahl der Singularitäten: Kapitel 12.}
\end{minipage}
\end{center}

\begin{center}
\begin{minipage}{0.88\textwidth}
\kursiv{\textbf{Jede rationale Kurve ist das Bild der Geraden, und das Bild ist meist singulär}: Die Gerade hat keine Knoten, aber die Abbildung kann zwei Werte $t = \pm1$ auf denselben Punkt werfen --- der Knoten --- oder in $t = 0$ die Ableitung null haben --- die Spitze. Die Singularitäten der rationalen Kurve sind die Stellen, an denen die Parametrisierung nicht injektiv oder nicht immersiv ist. Für den Leitfall ist die Hyperbel glatt und die Parametrisierung durch $k_{12}$ bijektiv (außer im Unendlichen); die vier Zeilen der Tabelle in Kapitel 1 sind vier Werte des Parameters.}
\end{minipage}
\end{center}
\end{merksatz}

\section{Lüroth}

\begin{merksatz}
\boxtitel{Satz 11.3 --- Lüroth 1876}
Sei $k\subsetneq K\subseteq k(t)$ ein Zwischenkörper des rationalen Funktionenkörpers. Dann ist $K = k(u)$ für eine rationale Funktion $u\in k(t)$. Geometrisch: Jede Kurve, die Bild der Geraden unter einer rationalen Abbildung ist, hat auch eine \kursiv{bijektive} rationale Parametrisierung (bis auf endlich viele Punkte).

\schritt{Beweisidee}
Sei $u = p(t)/q(t)\in K\setminus k$ mit teilerfremden $p, q$ und $n = \max(\deg p, \deg q)$ minimal unter allen Elementen von $K\setminus k$. Dann ist $t$ Nullstelle von $p(X) - uq(X)\in K[X]$, also $[k(t) : K]\le n$. Andererseits ist $[k(t) : k(u)] = n$ (das Polynom $p(X) - uq(X)$ ist irreduzibel über $k(u)$, Gauß-Lemma), und $k(u)\subseteq K$; also $[k(t) : K]\le n = [k(t) : k(u)]$ zwingt $K = k(u)$ (Auslassung 16, \zit{10}, Kap.~IV.2; der Kern ist die Irreduzibilität). $\blacksquare$

\schritt{Was das für Parametrisierungen heißt}
Die Parametrisierung $t\mapsto(t^{2}, t^{4})$ der Parabel $y = x^{2}$ ist zwei-zu-eins; Lüroth sagt: Es gibt eine ein-zu-eins, nämlich $u = t^{2}$, also $u\mapsto(u, u^{2})$. Allgemein lässt sich jede nicht-injektive Parametrisierung durch eine injektive ersetzen --- die \Begriff{eigentliche} Parametrisierung ---, und der Grad der Abbildung ist der Index $[k(t) : K]$.

\begin{center}
\begin{minipage}{0.88\textwidth}
\kursiv{\textbf{Lüroth ist der Satz, der \enquote{rational} zu einem Begriff macht}: Es kommt nicht darauf an, \kursiv{wie} man parametrisiert, sondern nur, ob der Funktionenkörper der Kurve ein $k(u)$ ist --- und die Dimension-eins-Fassung des Problems, ob jede \enquote{unirationale} Varietät rational ist. Für Flächen gilt das über $\mathbb C$ noch (Castelnuovo 1894), über anderen Körpern nicht (Zariski 1958); ab Dimension drei ist es falsch (Clemens--Griffiths, Iskovskikh--Manin, Artin--Mumford, alle 1972): Es gibt dreidimensionale Varietäten, die Bild des $\mathbb P^{3}$ sind, ohne zu ihm birational zu sein --- die kubische Dreifaltigkeit ist das berühmteste Beispiel, und ihr Beweis braucht Hodge-Theorie, Kapitel 15.}
\end{minipage}
\end{center}
\end{merksatz}

\section{Der Bruchfall}

\begin{warnung}
\boxtitelr{Rational über $\mathbb C$, nicht über $\mathbb Q$ --- und gar nicht}
\begin{center}
{\small
\renewcommand{\arraystretch}{1.5}
\begin{tabular}{@{}P{2.8cm}P{2.8cm}P{2.8cm}@{}}
\toprule
\textbf{Kurve} & \textbf{über $\mathbb C$} & \textbf{über $\mathbb Q$}\\
\midrule
$x^{2} + y^{2} = 1$ & rational & rational\\
$x^{2} + y^{2} = 3$ & rational & keine Punkte, nicht rational\\
$x^{2} + y^{2} = -1$ & rational & leer; über $\mathbb R$ leer\\
$y^{2} = x^{3} - 2x$ & nicht rational & nicht rational\\
$x^{3} + y^{3} = 1$ & nicht rational & nicht rational --- Fermat für $n = 3$\\
\bottomrule
\end{tabular}}
\end{center}

\begin{center}
\begin{minipage}{0.9\textwidth}
\textbf{Rationalität hat zwei Hindernisse: ein geometrisches --- das Geschlecht, über $\mathbb C$ sichtbar --- und ein arithmetisches --- das Fehlen eines Punktes über dem Grundkörper. Das zweite verschwindet über $\mathbb C$, das erste nie.}
\end{minipage}
\end{center}

\begin{center}
\begin{minipage}{0.88\textwidth}
\kursiv{\textbf{Dass die glatte Kubik nicht rational ist, hat dieses Kapitel noch nicht bewiesen}, nur behauptet. Der Beweis über die Gruppe ($\mathbb Q$ ist nicht endlich erzeugt, die Punkte einer rationalen Kurve wären es aber --- nein: das zeigt nur, dass sie nicht über $\mathbb Q$ rational ist) reicht nicht; über $\mathbb C$ braucht man eine Invariante, die die Gerade von der Kubik unterscheidet. Kapitel 12 gibt sie: das Geschlecht, berechenbar aus dem Grad und den Singularitäten, $0$ für die Gerade, $1$ für die glatte Kubik. Und für $x^{3} + y^{3} = 1$: Geschlecht $1$, also eine elliptische Kurve --- und ihre rationalen Punkte sind nur $(1, 0)$, $(0, 1)$ und der Punkt im Unendlichen, die Torsionsgruppe $\mathbb Z/3$, Rang $0$: Eulers Beweis von Fermats Satz für $n = 3$, in heutiger Sprache.}
\end{minipage}
\end{center}
\end{warnung}

\section{Übungsaufgaben}

\begin{aufgabe}
\boxtitel{Aufgabe 11.1 --- Das Folium}
Parametrisieren Sie $x^{3} + y^{3} = 3xy$ und finden Sie die Punkte im Unendlichen.

\loesung
\schritt{Schritt 1 --- Der Ursprung}
$f = x^{3} + y^{3} - 3xy$: Terme niedrigsten Grades $-3xy$, Vielfachheit $2$ (Knoten mit Tangenten $x = 0$, $y = 0$); Grad $3$. Satz 11.2 greift.

\schritt{Schritt 2 --- Gerade $y = tx$}
$x^{3} + t^{3}x^{3} - 3tx^{2} = 0$: $x = 0$ (doppelt) oder $x(1 + t^{3}) = 3t$: $x = \frac{3t}{1 + t^{3}}$, $y = \frac{3t^{2}}{1 + t^{3}}$.

\schritt{Schritt 3 --- Probe}
$t = 1$: $(\tfrac32, \tfrac32)$: $\tfrac{27}{8} + \tfrac{27}{8} = \tfrac{27}{4} = 3\cdot\tfrac94$. \checkmark\ $t = 2$: $(\tfrac69, \tfrac{12}{9}) = (\tfrac23, \tfrac43)$: $\tfrac{8}{27} + \tfrac{64}{27} = \tfrac{72}{27} = \tfrac83 = 3\cdot\tfrac89$. \checkmark

\schritt{Schritt 4 --- Unendlich}
Homogen $x^{3} + y^{3} = 3xyz$; $z = 0$: $x^{3} + y^{3} = 0$, $y = -x$ (reell), $y = -\omega x$, $y = -\omega^{2}x$ (komplex, $\omega^{3} = 1$). Ein reeller Punkt $(1 : -1 : 0)$, die Richtung der Asymptote $x + y = -1$; in der Parametrisierung bei $t = -1$ ($1 + t^{3} = 0$). Die zwei komplexen bei $t = -\omega, -\omega^{2}$.

\begin{center}
\begin{minipage}{0.88\textwidth}
\kursiv{\textbf{Das Folium ist Descartes' Kurve von 1638}, mit der er Fermats Tangentenmethode herausforderte --- Fermat löste die Aufgabe, Descartes nicht; die Asymptote fand erst \Person{Huygens} 1692. Die Parametrisierung macht alles leicht: Tangente in $t$ durch Ableiten nach $t$, Asymptote bei $t\to-1$. Was die Geometrie des 17.\ Jahrhunderts mühsam fand, ist für eine rationale Kurve eine Rechnung in einer Veränderlichen --- das ist der praktische Wert der Rationalität, und der Grund, warum die Konstrukteure von Kurven am Rechner (Bézier, NURBS) nur rationale Kurven verwenden.}
\end{minipage}
\end{center}
\end{aufgabe}

\begin{aufgabe}
\boxtitel{Aufgabe 11.2 --- Pythagoreische Tripel, vollständig}
Zeigen Sie mit der Parametrisierung des Kreises, dass jedes primitive pythagoreische Tripel $(a, b, c)$ mit $b$ gerade die Form $(m^{2} - n^{2},\ 2mn,\ m^{2} + n^{2})$ hat.

\loesung
\schritt{Schritt 1 --- Vom Tripel zum Punkt}
$(a/c, b/c)$ ist ein rationaler Punkt auf $x^{2} + y^{2} = 1$, verschieden von $(-1, 0)$.

\schritt{Schritt 2 --- Parametrisierung durch $(-1, 0)$}
Gerade $y = t(x + 1)$: $x^{2} + t^{2}(x + 1)^{2} = 1$, zweite Nullstelle $x = \frac{1 - t^{2}}{1 + t^{2}}$, $y = \frac{2t}{1 + t^{2}}$ (Kapitel 1). Der Punkt $(a/c, b/c)$ hat $t = \frac{b/c}{a/c + 1} = \frac{b}{a + c}\in\mathbb Q$.

\schritt{Schritt 3 --- $t = n/m$ gekürzt}
$x = \frac{m^{2} - n^{2}}{m^{2} + n^{2}}$, $y = \frac{2mn}{m^{2} + n^{2}}$. Also $(a, b, c) = \lambda(m^{2} - n^{2}, 2mn, m^{2} + n^{2})$ für ein rationales $\lambda$.

\schritt{Schritt 4 --- $\lambda = 1$}
Sind $m, n$ teilerfremd und von verschiedener Parität, so ist $(m^{2} - n^{2}, 2mn, m^{2} + n^{2})$ primitiv (ein gemeinsamer Primteiler $p$ teilte $2m^{2}$ und $2n^{2}$, also $p = 2$; aber $m^{2} + n^{2}$ ist ungerade). Zwei primitive Tripel, die proportional sind, sind gleich: $\lambda = 1$. Sind $m, n$ beide ungerade, so ist das Tripel durch $2$ teilbar, und $\tfrac12(m^{2} - n^{2}, 2mn, m^{2} + n^{2})$ hat die gerade Komponente an erster Stelle --- der Fall \enquote{$a$ gerade}, ausgeschlossen. $\blacksquare$

\begin{center}
\begin{minipage}{0.88\textwidth}
\kursiv{\textbf{Euklid X.29 in einer Rechnung}: Die vollständige Beschreibung aller pythagoreischen Tripel, die Euklid auf zwei Seiten führt, ist die Aussage, dass der Kreis rational ist, plus Kürzen. Das Muster überträgt sich auf jeden Kegelschnitt mit einem rationalen Punkt: $x^{2} - 2y^{2} = 1$ (Pell) hat die Parametrisierung durch $(1, 0)$, und ihre ganzzahligen Punkte sind die Lösungen der Pellschen Gleichung, Band 1 --- aber die Ganzzahligkeit ist dann eine zusätzliche Bedingung an $t$, die die Parametrisierung allein nicht liefert.}
\end{minipage}
\end{center}
\end{aufgabe}

\begin{aufgabe}
\boxtitel{Aufgabe 11.3 --- Rationale Funktionen auf dem Leitfall}
Die identifizierbaren Funktionen des Leitfalls sind konstant auf der Kurve $C$. Was sind die \kursiv{nicht} konstanten rationalen Funktionen auf $C$, und was hat das mit Lüroth zu tun?

\loesung
\schritt{Schritt 1 --- Funktionenkörper}
$k(C)$ ist der Quotientenkörper von $k[C] = k[k_{12}, k_{20}]/\langle g\rangle$ (Kapitel 3). Da $g$ linear in $k_{20}$ ist, lässt sich $k_{20}$ eliminieren: $k_{20} = \tfrac25 - \frac{3}{50k_{12}}$. Also $k(C) = k(k_{12})$: der rationale Funktionenkörper in einer Veränderlichen.

\schritt{Schritt 2 --- Die Funktionen}
Jede rationale Funktion auf $C$ ist eine rationale Funktion von $k_{12}$ allein: $k_{10} = \tfrac12 - k_{12}$, $k_{21} = \frac{3}{50k_{12}}$, $k_{20} = \tfrac25 - \frac{3}{50k_{12}}$, $k_{10}k_{21} = \frac{3(1 - 2k_{12})}{100k_{12}}$, \ldots\ Die konstanten sind genau die identifizierbaren; alle anderen hängen vom freien Parameter ab.

\schritt{Schritt 3 --- Lüroth}
Jeder Zwischenkörper $k\subsetneq K\subseteq k(k_{12})$ ist von der Form $k(u)$: Wählt man \kursiv{eine} nicht identifizierbare Funktion $u$ als zusätzlich gemessen --- etwa $u = k_{21}$ durch eine Gewebeprobe ---, so wird alles identifizierbar, was in $k(u)$ liegt. Für $u = k_{21} = \frac{3}{50k_{12}}$ ist $k(u) = k(k_{12})$: eine Messung genügt für alles. Für $u = k_{12}^{2}$ wäre $[k(k_{12}) : k(u)] = 2$: Alles wäre bis auf ein Vorzeichen bestimmt --- die $2$ aus Aufgabe 1.3 hat hier ihr Gegenstück.

\begin{center}
\begin{minipage}{0.88\textwidth}
\kursiv{\textbf{Die Frage \enquote{welche zusätzliche Messung macht das Modell identifizierbar?} ist eine Frage nach dem Funktionenkörper der Faser}, und für rationale Fasern beantwortet Lüroth sie vollständig: Jede Messung, die eine nichtkonstante Funktion liefert, erzeugt einen Zwischenkörper $k(u)$, und der Grad $[k(C) : k(u)]$ ist die Zahl der Modelle, die danach noch ununterscheidbar sind. Für Fasern, die elliptische Kurven sind, gibt es kein $u$ mit $k(C) = k(u)$ --- dann macht keine einzelne rationale Messgröße alles identifizierbar, und man braucht zwei.}
\end{minipage}
\end{center}
\end{aufgabe}

\begin{geschichte}
\boxtitelg{Diophant, Euklid, Lüroth}
Die Parametrisierung des Kreises steht bei \Person{Euklid} (Elemente X.29, Lemma 1) als Konstruktion aller pythagoreischen Tripel und bei \Person{Diophant} (Arithmetica II.8) als Methode, eine Quadratzahl in zwei Quadrate zu zerlegen --- die Stelle, an deren Rand \Person{Fermat} 1637 seine Bemerkung schrieb.

\begin{center}
\begin{minipage}{0.9\textwidth}
\kursiv{\Person{Jakob Lüroth}, Professor in Karlsruhe und Freiburg, bewies seinen Satz 1876 in den Mathematischen Annalen; \Person{Clebsch} hatte ihn 1868 für den Spezialfall der ebenen Kurven gezeigt, und \Person{Steinitz} gab 1910 den rein algebraischen Beweis. \textbf{Dass das Problem in höherer Dimension offen war, galt bis 1972 als eines der Hauptprobleme der Geometrie}; die drei Gegenbeispiele erschienen im selben Jahr mit drei völlig verschiedenen Methoden --- Hodge-Theorie (Clemens--Griffiths), birationale Starrheit (Iskovskikh--Manin), Torsion in der Kohomologie (Artin--Mumford) --- und begründeten die moderne birationale Geometrie.}
\end{minipage}
\end{center}

\schritt{Kurven am Rechner}
\begin{center}
\begin{minipage}{0.9\textwidth}
\kursiv{\Person{Pierre Bézier} bei Renault und \Person{Paul de Casteljau} bei Citroën entwickelten um 1960 die parametrischen Kurven, mit denen seither Karosserien, Schriften und jede Zeichnung am Rechner beschrieben werden --- polynomiale Parametrisierungen, später rationale (NURBS). \textbf{Rationalität ist hier keine Eigenschaft, die man nachweist, sondern eine, die man voraussetzt}: Ein Kreis ist in einem Zeichenprogramm die rationale Parametrisierung aus Kapitel 1, in Stücke geschnitten, und eine elliptische Kurve kann man dort nur näherungsweise zeichnen.}
\end{minipage}
\end{center}
\end{geschichte}

\begin{forschung}
\boxtitelp{Rationalität heute}
\begin{center}
{\small
\renewcommand{\arraystretch}{1.45}
\begin{tabular}{@{}P{3.0cm}P{5.2cm}@{}}
\toprule
\textbf{Rationalität von Hyperflächen} & Welche glatten Hyperflächen vom Grad $d$ im $\mathbb P^{n}$ sind rational? Für kubische Vierfaltigkeiten im $\mathbb P^{5}$ offen und heftig umstritten --- die Vermutung: die allgemeine nicht, spezielle ja (Hassett 2000, Kuznetsov 2010).\\
\textbf{Stabile Rationalität} & Voisin 2015: Die Spezialisierungsmethode zeigt, dass sehr allgemeine Quartik-Dreifaltigkeiten nicht einmal stabil rational sind; Schreieder 2019 für Hyperflächen fast aller Grade.\\
\textbf{Algorithmische Parametrisierung} & Für Kurven vom Geschlecht null: Parametrisierung über $\mathbb Q$ finden oder beweisen, dass keine existiert (van Hoeij 1997, Sendra--Winkler); implementiert in Maple und Magma.\\
\bottomrule
\end{tabular}}
\end{center}

\schritt{Zur dritten Zeile}
\begin{center}
\begin{minipage}{0.88\textwidth}
\kursiv{\textbf{Der Algorithmus zur Parametrisierung ist Kapitel 11 als Programm}: Geschlecht berechnen (Kapitel 12); ist es null, einen Punkt über dem Grundkörper suchen --- das ist der schwere Teil, Legendres Kriterium für Kegelschnitte, im Allgemeinen über einen adjungierten Kegelschnitt ---, dann Geraden durch den Punkt. \Person{Sendra} und \Person{Winkler}s Buch von 2008 ist die Referenz, und das Verfahren ist das, mit dem ein Computeralgebrasystem auf \enquote{parametrize} antwortet. Für den Leitfall tut es, was Kapitel 1 von Hand getan hat; für eine Faser vom Grad $6$ in zwei Veränderlichen tut es, was von Hand niemand täte.}
\end{minipage}
\end{center}
\end{forschung}

\begin{merksatz}
\boxtitel{Was Sie jetzt können}
\begin{center}
\renewcommand{\arraystretch}{1.5}
\begin{tabular}{@{}P{3.0cm}P{5.6cm}@{}}
\toprule
\textbf{Satz 11.1} & Kegelschnitt mit Punkt: Gerade durch den Punkt, Vieta\\
\textbf{Satz 11.2} & Vielfachheit $d - 1$: Geraden durch den singulären Punkt\\
\textbf{Satz 11.3} & Lüroth: Zwischenkörper von $k(t)$ sind $k(u)$\\
\textbf{Der Bruchfall} & zwei Hindernisse: Geschlecht (geometrisch) und Punkt (arithmetisch)\\
\bottomrule
\end{tabular}
\end{center}

\begin{center}
\begin{minipage}{0.88\textwidth}
\kursiv{Rational oder nicht --- die Antwort steht noch aus. Kapitel 12 gibt die Zahl, die sie entscheidet: das Geschlecht, aus Grad und Singularitäten berechnet, $0$ für alles, was dieses Kapitel parametrisiert hat, $1$ für die Kubik des Titelbilds.}
\end{minipage}
\end{center}
\end{merksatz}
